
%---------------- Conclusion ----------------------------------


\section{Conclusion}

In this paper, we have constructed an infinite family of invariants of algebraic curves.
By construction, these invariants coincide with the topological expansion of matrix integrals in the special case where the algebraic curve is the large $N$ part of the matrix integral's spectral curve.
But we emphasize again that the construction presented here goes beyond matrix models.

Our invariants are defined only in terms of algebraic geometry, and they have many interesting properties, like homogeneity, and integrability (they obey some Hirota equation).

The problem of computing the $F^{(g)}$'s for matrix models is an old problem which was addressed many times, and which found more and more elaborate answers \cite{ACKM}.
We claim that ours is more efficient, because it contains all multicut cases, and various types of matrix models at once. Also, even in the simplest cases (1-matrix model, 1 cut), our expressions are simpler than what existed before \cite{ACKM}.
Our $F^{(g)}$'s are defined recursively, like those of \cite{ACKM}, but the recursions are much easier to handle, and it is much easier to deduce properties to any order $g$ from our construction.

The efficiency of our method becomes striking when one wants to compare different models (Kontsevitch and KdV for instance), or when one wants to take singular limits.
Another important application of our method, is to prove that our $F^{(g)}$'s provide a solution to the holomorphic anomaly equations of \cite{BCOV} in topological string theory, thus confirming the Dijkgraaf Vafa correspondence. We claim that this can be proved easily from our work, and we present it in a coming paper \cite{EMO}.
It would also be interesting to compare our free energies with the $D$-Modules considered in \cite{Mth1,Mth2} as partition functions of a unified matrix M-theory by checking
that they indeed satisfy the equations of \cite{Mth2}.

One of the reasons our method is very efficient also, is that it can be represented diagrammatically, without equations, and thus very easy to remember.



\bigskip

{\bf Perspectives and generalizations:}


\begin{itemize}

\item The first thing one could think about is to understand what our $F^{(g)}$'s compute in algebraic geometry. There has been some attempts to recognize the first few of them $F^{(1)}$ as the Dedekind function \cite{Dubrovin1, Dubrovin2}, $F^{(2)}$ as the Eisenstein series \cite{klemm}, but the answer for higher $g$ is still obscure.
Also a combinatoric interpretation is missing.


\item Beyond that, it would be important to understand the combinatorics behind our diagrammatic construction. Since all diagrams are obtained from trees, it seems to be related to Schaeffer's method for counting maps \cite{schaeffer}, but this issue needs to be investigated further.

\item Double scaling limits of matrix models are in relationship with conformal field theory (CFT), and in particular minimal models. It would be interesting to compare our formulae with those obtained directely from CFT \cite{BookPDF, Zamozamo, Zamo}, and, since higher genus CFT is far less known, maybe our method can bring something new to CFT.

\item Also, the link between formal matrix models (i.e. those defined as combinatorial generating functions, for which it makes sense to consider a topological expansion, cf. \cite{eynform}), and actual convergent matrix integrals needs to be better understood. The difficulty lies beyond any order in perturbation, and integrability could play a key role in understanding the relationship in more details.

\item It would be interesting to check that the topological expansion for the chain of matrices \cite{eynmultimat},  is also given by the same $F^{(g)}$'s. This is an open question, but there are strong evidences that the answer is positive. For instance it is easy to see from \cite{eynmultimat} that $F^{(0)}$ and some of the first few correlation functions are indeed the same.

\item It would be interesting also to extend our construction to other types of matrix models, for instance non-hermitian (real symmetric or quaternionic). A first attempt was done in \cite{ChEynbeta}.
Another possible extension is towards the $O(n)$ matrix model, whose large $N$ limit is known in terms of an agebraic curve \cite{ECK1,ECK2}. To leading order, the $O(n)$ matrix model solution looks very similar to that of the 1-matrix model, except that correlation functions are no longer meromorphic functions on the curve. Instead they gain a phase shift after going around a non-trivial cycle. This would probably allow to define a twisted version of our construction.

\end{itemize}



\subsection*{Acknowledgements}

We would like to thank Mina Aganagic, Michel Berg\`ere, Leonid Chekhov, Philippe Di Francesco, Robert Dijkgraaf, Aleix Prats Ferrer, Ivan Kostov, Amir Kashani-Poor, Marcos Mari\~no, Andrei Okounkov, Pierre Vanhove, Jean-Bernard Zuber, for fruitful discussions on this subject.
This work is partly supported by the Enigma European network MRT-CT-2004-5652, by the ANR project G\'eom\'etrie et int\'egrabilit\'e en physique math\'ematique ANR-05-BLAN-0029-01,
 by the Enrage European network MRTN-CT-2004-005616, 
 by the European Science foundation through the Misgam program,
 by the French and Japaneese governments through PAI Sakurav, by the Quebec government with the FQRNT.













